\newpage
\section{Related Work}

Chapter 2 established what the evidence indicates should be effective. This chapter examines what has been built.

\subsection{Comparison, gap analysis, and positioning}

\subsubsection{Comparative assessment}

Ratings use three levels. Full ($\bullet$) indicates the attribute is a core designed property of the system. Partial ($\ominus$) indicates it is present but limited. Absent ($\circ$) indicates it is not provided.

\begin{landscape}
\begin{scriptsize}
\setlength{\tabcolsep}{3pt}
\renewcommand{\arraystretch}{1.3}
\begin{longtable}{@{}p{3.1cm} *{3}{p{5.9cm}}@{}}
\caption{Comparison of system classes against the derived criteria.}\\
\toprule
\textbf{Attribute} & \textbf{Authored simulation (Forage)} & \textbf{Generated simulation (EntryLevel)} & \textbf{Proposed system} \\
\midrule
\endfirsthead
\toprule
\textbf{Attribute} & \textbf{Authored simulation (Forage)} & \textbf{Generated simulation (EntryLevel)} & \textbf{Proposed system} \\
\midrule
\endhead
C1 \newline Task authenticity & {\small$\bullet$}\  employer-authored real work products & {\small$\bullet$}\  generated role tasks & {\small$\bullet$}\  tasks drawn from occupational taxonomies \\
C3 \newline Adaptivity to the learner & {\small$\circ$}\  fixed content, identical on repeat & {\small$\bullet$}\  parameterised generation & {\small$\bullet$}\  seeded generation per user \\
C4 \newline Validity of competence measurement & {\small$\circ$}\  model answer, no judgment & {\small$\ominus$}\  behavioural evidence, unvalidated & {\small$\ominus$}\  designed; validity to be established \\
\bottomrule
\end{longtable}
\end{scriptsize}

\noindent\scriptsize{$\bullet$~Full --- a core designed property \quad $\ominus$~Partial --- limited, incidental, or a by-product \quad $\circ$~Absent.}
\end{landscape}
